\begin{abstract}
Limited defensive budgets motivate prioritizing npm maintainers by downstream reach. However, reaching many repositories and reducing exposure to possible account compromises are different objectives. We study this distinction in our collected graph of \StudyPackages{} package nodes and \StudyMaintainers{} maintainer usernames. We compare five deterministic selection methods and 200 random orders across six budgets, using package and workflow-associated repository endpoints. At 100 identities, marginal coverage greedy reaches \GreedyCovered{} of \StudyWorkflowRepos{} observed workflow repositories (\GreedyCoveragePercent{}\%), versus \UnionCovered{} for exact maintainer-union ranking. Under a stipulated uniform single-account compromise model with perfect source hardening, their exposure reductions are \GreedyExposurePercent{}\% and \UnionExposurePercent{}\%, respectively. Only \GreedyExcluded{} repositories covered by greedy satisfy exclusion of every observed source identity. Preserving version continuity reduces identity--repository exposure pairs by \VersionPairsDecreasePercent{}\%, although aggregate greedy coverage remains similar. These results extend established maintainer-reach analyses with an explicit comparison of objectives and observation units. They support objective-specific prioritization and qualified structural measurements; they do not estimate incident probabilities or demonstrated control effectiveness.
\end{abstract}
